\documentclass[11pt]{article}
\usepackage{amsmath,amssymb,amsthm,hyperref}
\newtheorem{proposition}{Proposition}
\newtheorem{theorem}{Theorem}
\newtheorem{problem}{Problem}
\newcommand{\one}{\mathbf 1}
\newcommand{\R}{\mathbb R}
\newcommand{\F}{\mathbb F}
\title{Embedded equitable community modes:\\exact certificates and limits of Bethe--Hessian outliers}
\author{Sparse partitioning research session}
\date{15 November 2026}
\begin{document}
\maketitle

\paragraph{Status.}
The propositions below are elementary deterministic facts and explicit
counterexamples, not a claimed breakthrough or a new recovery threshold.
They isolate an exact random-matrix problem that would imply a substantial
extension of regular stochastic block model recovery. That probabilistic
problem is NOT proved here. The conjectural direction comes from Barucca
[Ba] and is explicitly described as open in the September 2026 revision
of Zhu--Zhu [ZZ, Section 1.3]. We inspected those actual primary sources.

Let $a>b>0$ be integers, $d=a+b$, $\lambda=a-b$, and $N$ be even.
A balanced $(a,b)$-equitable bisection of a simple graph is a sign vector
$\sigma\in\{-1,1\}^N$ with $\one^T\sigma=0$ such that every vertex
has exactly $a$ neighbors of its own sign and $b$ of the opposite sign.
Its two signs describe an unordered partition, so $\sigma$ and $-\sigma$
are equivalent. Write $A$ for adjacency and
\[
 H(r)=(r^2-1)I+D-rA
 \quad(r\in\R)
\]
for the standard Bethe--Hessian convention. Other papers use the reciprocal
parameter $t$, namely $I-tA+t^2(D-I)=t^2H(1/t)$ when $t\ne0$.

\begin{proposition}[Exact rank certificate]
Suppose $G$ is $d$-regular and has a balanced $(a,b)$-equitable bisection.
Then $A\sigma=\lambda\sigma$. If
\[
 \dim_{\R}\ker(A-\lambda I)=1,                 \tag{1}
\]
then the bisection is unique up to sign and can be recovered by exact
linear algebra in polynomial bit complexity. Conversely, any balanced
sign vector satisfying $A\tau=\lambda\tau$ defines an $(a,b)$-equitable
bisection. No converse from uniqueness of sign bisections to (1) is claimed.
\end{proposition}
\begin{proof}
For vertex $v$, the signed neighbor sum is $(a-b)\sigma_v$.
Conversely, if $x_v$ denotes the number of opposite-sign neighbors of a
sign eigenvector, then $(d-2x_v)\tau_v=\lambda\tau_v$, whence
$x_v=(d-\lambda)/2=b$. Under (1), every sign eigenvector is a scalar
multiple of $\sigma$ and thus equals $\pm\sigma$. Gaussian elimination
over $\mathbb Q$, with fraction-free elimination for example, finds this
one-dimensional kernel in polynomial bit complexity; normalization and
exact verification recover the signs. No separation from nearby real
eigenvalues is required for this exact arithmetic statement.
\end{proof}

\begin{proposition}[A single cofactor and signed spanning-tree cancellation]
Assume the same equitable-bisection promise and put
$P_A(t)=\det(tI-A)$. For a spanning tree $T$ of $G$, let $c_\sigma(T)$
be the number of its edges crossing the bisection, and set
\[
 Z_\sigma=\sum_{T\text{ spanning tree}}(-1)^{c_\sigma(T)}.
\]
Then every principal $(N-1)$-cofactor of $\lambda I-A$ equals
$Z_\sigma$, and
\[
 P_A'(\lambda)=N Z_\sigma.                         \tag{6}
\]
In particular, (1) is equivalent to $Z_\sigma\ne0$, and to nonvanishing
of ANY ONE principal $(N-1)$-minor of $A-\lambda I$. The latter test
can be performed without knowing $\sigma$. This is the weighted
matrix-tree theorem specialized to the planted eigenvector, not a new
matrix-tree identity.
\end{proposition}
\begin{proof}
Let $S=\operatorname{diag}(\sigma)$ and $C=SAS$. Every row sum of $C$
equals $\lambda$ because $A\sigma=\lambda\sigma$. Therefore
\[
 L=\lambda I-C=S(\lambda I-A)S
\]
is the weighted Laplacian whose edge $uv$ has weight
$w_{uv}=\sigma_u\sigma_v$, namely $+1$ inside a community and $-1$
across it. If $R$ is any oriented vertex-edge incidence matrix, then
$L=R\operatorname{diag}(w)R^T$. Delete row $v$ from $R$. The
Cauchy--Binet expansion of the resulting principal Laplacian determinant
sums $\det(R_{\widehat v,T})^2\prod_{e\in T}w_e$ over all edge sets
of size $N-1$. The incidence determinant is zero unless the edge set
is a spanning tree, and equals $\pm1$ for a spanning tree. This proves
that the cofactor equals $Z_\sigma$, independently of $v$. Conjugation
by $S$ does not change any principal minor. Differentiating the
determinant in $t$ sums the $N$ principal cofactors, proving (6).

Because $\lambda I-A$ is real symmetric, $P_A'(\lambda)\ne0$ holds
exactly when its zero eigenvalue has multiplicity one. Equivalently all
the displayed principal cofactors are nonzero. Multiplying the minor by
$(-1)^{N-1}$ transfers the conclusion to $A-\lambda I$.
\end{proof}

\begin{proposition}[A checkable modular version]
An explicit balanced sign vector $\sigma$ with $A\sigma=\lambda\sigma$,
together with a prime $p$ and one principal $(N-1)\times(N-1)$ minor of
$A-\lambda I$ nonzero modulo $p$, certifies (1) and unique recovery.
Equivalently, it suffices to verify
$\operatorname{rank}_{\F_p}(A-\lambda I)=N-1$.
For a graph promised to possess the bisection, an odd prime of modular
nullity one permits its recovery directly from the modular kernel,
followed by exact integer verification.
\end{proposition}
\begin{proof}
A nonzero minor modulo $p$ is a nonzero integer minor, so the rational
rank is at least $N-1$; the verified integer kernel gives the reverse
inequality. Under the promise, every nonzero vector in the modular kernel
has the form $c\sigma$ in $\F_p^N$. Since $p$ is odd, its two values
$c,-c$ are distinct and determine the signs, which one then verifies in
integers. Modular rank failure is inconclusive. In particular if
$p\mid d-\lambda=2b$ with $p$ odd, both $\one$ and $\sigma$ belong
to the modular kernel and are independent; such a prime must fail.
Even if the rational rank is $N-1$, a fixed small prime can divide every
maximal minor. There is no assertion that a fixed-prime test succeeds
with probability tending to one.
\end{proof}

\begin{proposition}[Sharp scalar Bethe--Hessian barrier]
For every equitable bisection of a $d$-regular graph,
\[
 H(r)\sigma=h_\lambda(r)\sigma,
 \qquad h_\lambda(r)=r^2-\lambda r+d-1
 =\left(r-\lambda/2\right)^2+d-1-\lambda^2/4.       \tag{2}
\]
Consequently, if $\lambda^2<4(d-1)$, the signal direction is orthogonal
to the negative spectral subspace of $H(r)$ for EVERY real $r$.
If $\lambda^2=4(d-1)$ it is never negative and only touches zero at
$r=\lambda/2$. If $\lambda^2>4(d-1)$ it is negative exactly between
the real roots
\[
 r_\pm=\frac{\lambda\pm\sqrt{\lambda^2-4(d-1)}}2.   \tag{3}
\]
Thus the condition for some scalar Bethe--Hessian to have a negative
eigenvalue in the exact planted direction is precisely strict adjacency
outlier separation, $|\lambda|>2\sqrt{d-1}$.
\end{proposition}
\begin{proof}
Use $D=dI$ and $A\sigma=\lambda\sigma$ and complete the square.
Distinct real eigenvalues of the symmetric matrix $H(r)$ have orthogonal
eigenspaces, which proves the orthogonality assertion. The quadratic
has a negative value exactly when its discriminant is positive.
\end{proof}

The barrier concerns algorithms that retain only negative eigenspaces or
count negative modes. It does not rule out algorithms that inspect an
interior positive eigenvector, use exact algebraic knowledge of $\lambda$,
or process the graph with a different operator. Taking all real parameters
$r$ does not remedy this particular negative-eigenspace barrier.

\begin{proposition}[The corresponding NB modes lie on the bulk circle]
Let $d\geq3$ and $\lambda^2<4(d-1)$. Set
\[
 \theta_\pm=\frac{\lambda\pm i\sqrt{4(d-1)-\lambda^2}}2.
\]
Then $|\theta_\pm|=\sqrt{d-1}$, and the nonbacktracking operator has
eigenvectors
\[
 f_\theta(u,v)=\sigma_v-\theta^{-1}\sigma_u
 \quad\hbox{on directed edges }(u,v)
\]
with eigenvalues $\theta=\theta_\pm$. These are nonreal modes exactly
on the regular-graph bulk circle. They are NOT the real ``insider'' modes
of the separated regime in [ZZ].
\end{proposition}
\begin{proof}
The roots satisfy $\theta^2-\lambda\theta+d-1=0$ and are conjugates.
For the convention $(Bf)(u,v)=\sum_{w\sim v,w\ne u}f(v,w)$,
\[
 (Bf_\theta)(u,v)
 =\lambda\sigma_v-\sigma_u-(d-1)\theta^{-1}\sigma_v
 =\theta\sigma_v-\sigma_u=\theta f_\theta(u,v).
\]
Nonreality of $\theta$ and the nonzero real signs ensure $f_\theta\ne0$.
The modulus follows from the root product $d-1$.
\end{proof}

\paragraph{First tests.}
For $(a,b)=(2,1)$, $d=3$, $\lambda=1$,
\[
 h_1(r)=(r-1/2)^2+7/4>0,\qquad
 \theta_\pm=(1\pm i\sqrt7)/2.
\]
This is the smallest assortative cubic test, but it lies outside the
standing $\min(a,b)\geq3$ assumption of the original rigorous RSBM paper
[BDGHT]. The smallest assortative pair within that assumption and with
the even-internal-degree hypothesis used in its random-lift comparison is
$(a,b)=(4,3)$. There $d=7$, $\lambda=1$,
$h_1(r)=(r-1/2)^2+23/4>0$ and
$\theta_\pm=(1\pm i\sqrt{23})/2$.

\begin{problem}[A substantial missing probabilistic input]
For fixed integers $a>b\geq3$ with $a$ even, let $A_N$ be adjacency
of the uniform simple two-community regular block model. Prove
\[
 \Pr\!\left\{\dim\ker\bigl(A_N-(a-b)I\bigr)=1\right\}
 \longrightarrow1.                                           \tag{4}
\]
The especially concrete embedded test is $(a,b)=(4,3)$. The analogous
claim for every fixed $a>b>0$, including $(2,1)$, would establish
Barucca's full exact-recovery conjecture by the first proposition.
No proof of (4) is supplied.
\end{problem}

For even $a$, the permutation analogue is a lift of the two-vertex
multigraph with $a/2$ loops at each vertex and $b$ edges between them.
The old sign mode has eigenvalue $a-b$. Statement (4) asks that the
new-space adjacency have no eigenvalue EXACTLY equal to this old value.
The known almost-Ramanujan theorem excludes it when $a-b>2\sqrt{d-1}$,
but cannot exclude an exact collision at an energy inside the universal
cover spectrum. Local laws at scale $\log^C N/N$ and delocalization
also do not, on their own, exclude one extra eigenvalue at a prescribed
energy: one such eigenvalue contributes only $O(\log^{-C}N)$ to a
normalized smoothed density at that scale.

\begin{proposition}[The exact certificate is brittle under an edge switch]
Suppose $uv$ and $xy$ are within-community edges with
$\sigma_u=\sigma_v=1$ and $\sigma_x=\sigma_y=-1$, and suppose $ux,vy$
are absent. Delete $uv,xy$ and insert $ux,vy$. This preserves degrees and
simplicity, but for its adjacency $A'$,
\[
 (A'-\lambda I)\sigma=-2\sigma\one_{\{u,v,x,y\}}.  \tag{5}
\]
In particular $\sigma$ is no longer an exact eigenvector, even though
\[
 \|(A'-\lambda I)(\sigma/\sqrt N)\|_2=4/\sqrt N\longrightarrow0.
\]
\end{proposition}
\begin{proof}
At each of the four endpoints a same-sign neighbor is replaced with an
opposite-sign neighbor, changing the signed neighbor sum by $-2\sigma_v$.
All other rows have unchanged signed neighbor sums.
\end{proof}

Thus (4) is an exact-model algebraic target. This proposition does not
claim that robust recovery itself is impossible, and it supplies no
robust replacement for the exact rank certificate.

\begin{theorem}[Flat spectral ambiguity even at an optimal cover gap]
For every $n$, every $n$-lift of the cube $Q_3$ has three mutually
orthogonal balanced $(2,1)$-equitable sign vectors. Their normalized
inverse participation ratios are all exactly $1/(8n)$, and their
Bethe--Hessian eigenvalue curves are identical and strictly positive for
all real parameters. A random $n$-lift is connected with probability
$1-O(n^{-4})$, is locally tree-like, and, after discarding the two
trivial bipartite eigenvalues $\pm3$, has spectrum inside
$[-2\sqrt2-o_{\Pr}(1),2\sqrt2+o_{\Pr}(1)]$.

If a hidden community assignment is chosen independently and uniformly
from these three unordered partitions, no algorithm observing the graph
can recover that designated assignment with probability greater than
$1/3$. This conclusion remains true if the graph is independently
relabelled before it is observed.
\end{theorem}
\begin{proof}
Index cube vertices by $x=(x_1,x_2,x_3)\in\{0,1\}^3$, with edges
flipping one bit, and write lift vertices $(x,i)$, $i\in[n]$. Let
$\sigma_j(x,i)=(-1)^{x_j}$. Each vertex has two neighbors preserving
$\sigma_j$ and one reversing it, giving $A\sigma_j=\sigma_j$.
Uniform fiber sizes give balance, orthogonality, and exactly flat
normalized magnitudes. Apply (2) at $d=3,\lambda=1$.

To see connectivity, gauge the seven edge permutations on a spanning
tree to the identity. The five chord permutations remain independent
and uniform in $S_n$. The lift is connected if and only if their action
is transitive. If not, a nonempty invariant subset of size at most $n/2$
exists, so a union bound gives
\[
 \Pr\{\hbox{disconnected}\}
 \leq\sum_{k=1}^{\lfloor n/2\rfloor}\binom nk^{1-5}
 =O(n^{-4}).
\]
For fixed radius $R$, exposing a lifted neighborhood encounters a
collision with probability $O(4^R/n)$; thus a uniformly sampled rooted
neighborhood converges to the cubic tree. The old cube spectrum is
$3,1,1,1,-1,-1,-1,-3$. The fixed-base random-lift almost-Ramanujan theorem
[Bo] puts every new eigenvalue within $2\sqrt2+o_{\Pr}(1)$ in absolute
value, which proves the full nontrivial spectral assertion.

Finally condition on the sampled lift and on any independent vertex
relabeling. An algorithm's output is then independent of the hidden
uniform index. It can equal at most one of the three distinct unordered
partitions, so its conditional success probability is at most $1/3$.
Randomized algorithms follow by conditioning on their random seed.
\end{proof}

The cube example is not a sample from the uniform RSBM. It shows why
regular degrees, full flatness, all scalar Bethe--Hessian parameter
curves, local tree likeness, and an optimal universal-cover spectral gap
are insufficient deterministic hypotheses for unique recovery in the
embedded regime. The genuine random-model collision exclusion in (4)
is additional information.

\begin{thebibliography}{9}
\bibitem{Ba} P. Barucca, \emph{Spectral partitioning in equitable graphs},
Phys. Rev. E 95 (2017), 062310.
\url{https://arxiv.org/abs/1610.02668}, inspected v5, pp. 4--5.
\bibitem{ZZ} X. Zhu and Y. Zhu, \emph{Non-backtracking eigenvalues and
eigenvectors of random regular graphs and hypergraphs}, v3 (25 Sep 2026).
\url{https://arxiv.org/html/2312.03300v3}.
Inspected Section 1.3, Definition 3.9, Theorem 3.10, Corollary 3.11,
Remark 3.12, Lemma 4.1 and Section 9.
\bibitem{BDGHT} G. Brito, I. Dumitriu, S. Ganguly, C. Hoffman, L.V. Tran,
\emph{Recovery and rigidity in a regular stochastic block model},
SODA 2016; arXiv v4 (4 May 2020),
\url{https://arxiv.org/abs/1507.00930}.
Theorem 2 and Sections 1, 2.3, 4 were inspected.
\bibitem{Bo} C. Bordenave, \emph{A new proof of Friedman's second
eigenvalue theorem and its extension to random lifts},
\url{https://arxiv.org/abs/1502.04482}.
This is one of the repository's exact problem-21 sources.
\bibitem{HX} J. Huang and R. Xu, \emph{Local law and eigenvector
delocalization for random lifts}, 27 Oct 2026,
\url{https://arxiv.org/html/2610.37463}.
Inspected abstract, Sections 1.1--1.2, Theorems 1.2--1.3.
The base here is fixed, simple and of minimum degree at least three;
the two-vertex loop multigraph above is outside that stated scope.
\end{thebibliography}
\end{document}
